% This ensures VS Code compiles the main file even when you're editing a chapter.
% !TeX root = ../main.tex

\chapter{Introduction}
\label{cap:intro}

\section{The initial problem}

In order to navigate the real world, autonomous drones need to estimate their position
with reference to the world, with a certain precision \cite{introLivro1}. When used in
rural areas, drones are often used to disperse seeds, pesticides, or survey the field.
Field surveillance and inspection can now be done remotely, without the need for the
farmer to traverse the entire plantation, which could be thousands of acres of land
 \cite{introArtigo1}. Urban environments also benefit from drone services, from preventive
 maintenance services \cite{droneurban1} to emergency response infrastructure \cite{droneurban2}
\cite{droneurban3}.

\gls{gnss} is the main positioning system used for drones \cite{introLivro2}. Although GNSS
systems allow for more precise flight procedures, they come with a drawback: More precise and
robust systems usually involve heavier and bulkier hardware. Absolute positioning may not be
necessary in most cases though: Sometimes, the position of the drone is not as critical as its
position relative to a ground target, sometimes called \gls{poi}. While the accuracy of
the absolute position of the drone can be increased over time through the use of filtering,
quick and precise relative positioning between the drone and the ground target is critical so
the drone can either approach it or avoid it.

The purpose of this dissertation is to investigate the feasibility of a sensor fusion system
for precise ground target localization, comprising of one distance-measuring pod, to be attached
to the drone as payload. The drone system itself is outside of the scope of this work.

\section{Mission scope}

As mentioned before, the drone absolute position acquisition system is not part of the scope of this
work. This dissertation will focus on the payload (also referred as pod) of the drone, to be connected
to the drone in order to measure distances. A first system will be developed to behave as a drone, and
through the use of computer vision will identify a test target and calculate its direction relative
to the drone. This direction information will be transmitted to the pod, as a set of angles, and the
pod will then measure the distance between the pod and target. It will also compensate for parallax, 
as the pod could be positioned anywhere on the drone, and then estimate the distance between the
camera and the object, as if the camera itself had measured it instead of an external module \cite{robotics}.
An example of the scenario can be seen in Figure~\ref{fig:1img1}. Concepts like what distance sensor
to use or \gls{fov} will be explained further on Chapter~\ref{cap:back}.

\begin{figure}[h]
  \centering
  \includegraphics[width=0.5\textwidth]{1/img1}
  \caption{Drone observing target on ground with a camera while measuring its distance}
  \label{fig:1img1}
\end{figure}

Different sensors have different specifications. In order to investigate the many devices available for
designing a precise drone location system, well-defined requirements are of the utmost importance. As this
survey was carried out in Brazil, all system specifications will be defined by
Brazilian regulations and plantation standards.

The first parameter to define is the altitude range in which the drone will fly. Brazil's Ministry of
Agriculture defines, as of November 2024 \cite{dronealturaagro}, that the maximum flight altitude allowed
for a drone used for agricultural purposes or far from crowded areas is 400 feet (approximately 120 meters).
In addition to this resolution, the \gls{anac} defines that any drone that flies above 400
feet must be registered and licensed \cite{dronealturaANAC}. This is, of course, the maximum allowed
operating altitude, and not necessarily the altitude at which the drone will, in fact, operate.

\section{System architecture}

This target location system will consist of a target-to-drone distance measuring system.
The pod should first receive orientation from the drone system. Then measure the
distance to the target on ground. To obtain the correct distance of the target to
the drone, the pod system should make some corrections necessary for the difference in point-of-view.


\section{Different technologies for target location}

Many different technologies can be used to detect the distance between a sensor and an object/obstacle
\cite{robotsens}.

\subsection{Ultrasonic sensors}

Ultrasonic sensors consist of a sensor that emits ultrasonic pulses (above 20 kHz)
and measures the time between the emission and the receiving of an echo \cite{ultrasound1}.
Ultrasonic sensors are commonly used in robotics, automotive parking systems, and level detection
because of their ability to detect objects without requiring physical contact. They are inexpensive,
work well under various lighting conditions, and are effective in detecting both solid and liquid
surfaces, although their accuracy can be affected by environmental factors such as temperature,
air pressure, and the type of surface to be detected. Ultrasonic sensors rarely measure distances
greater than 10 meters \cite{ultrasound2}.

\subsection{Light-based sensors}

Similarly to radio waves systems, there are equivalent sensors that utilize electromagnetic radiation
in the visible spectrum. The main types of light-based sensors are the following
 \cite{infrared}\cite{laserrangefinder}:

\begin{itemize}
  \item \gls{tof} : These sensors emit light (typically laser or \gls{ir}) and measure
  the time it takes for the reflected light to return to the sensor. They are widely used in
  smartphones, robotics, and drones because of their precision and ability to operate over a range
  of distances. These modules detect a single-point of measurement.
  \item \gls{lidar}: These use laser beams to measure distance by timing the
  reflected light. \gls{lidar} systems are commonly used in autonomous vehicles, mapping, and industrial
  applications due to their high accuracy and ability to create 3D maps. They usually consist of a light
  transmitter and receiver module rotating, in order to get wide angles of measurements, or \gls{tof}
  sensor that detect a grid or matrix of light beams, and not just a single point.
  \item Structured Light Sensors: These sensors project a known light pattern (for example, a grid)
  onto an object and analyze the deformation of the pattern to calculate distance. They are used in
  depth-sensing cameras like Microsoft Kinect for 3D scanning and motion tracking.
  \item Stereo Vision: This method uses two or more cameras to simulate human binocular vision.
  By analyzing the disparity between images captured by each camera, it calculates the depth and distance.
  Stereo vision is used in robotics, AR/VR, and object recognition systems.
  \item RGB-D/Depth Camera: Depth cameras are a combination of \gls{tof} sensors and standard cameras.
  The combination allows for the distance information to be obtained together with the image but it also
  combines the drawbacks of matrix \gls{tof} sensors, with added cost and processing complexity.
\end{itemize}

Structured light sensors are more commonly used in civil engineering, as they are usually too
bulky to be fitted on drones. Stereo vision is by far the most complex of all the methods cited,
being very hard to calibrate. \gls{tof} and \gls{lidar} are preferred in drone applications and other
autonomous robots for their simplicity, size, and wide range of specifications \cite{dronesensors}.

\subsubsection{Single point ToF vs. multi-point LiDAR}

A \gls{tof} sensor will measure the distance between the sensor light pulse transmitter and receiver.
As the transmitter and receiver are pointed in the same direction, an object must be place in front of
the sensor to bounce the light pulse transmitted back to the receiver. The measurement consists of measuring
the time this pulse takes to be detected by the receiver after the transmitter was activated, and the distance
is divided by two, for the light pulse traversed the distance measured twice \cite{lidar2}.

A \gls{lidar} sensor will make multipledistance measurements from different points in space \cite{lidar1}. This
can be done by moving a \gls{tof} sensor in a circular motion (sometimes called scanning \gls{lidar}) or using
a multipoint solid-state \gls{tof} sensor (sometimes called flash \gls{lidar}). Multipoint \gls{tof} sensors
consist on a \gls{tof} sensor that generates multiple light beams, in a line or in a matrix, in order to measure
distances from multiple points in a 3D space.

While \gls{lidar} sensors can be very practical, they are bulkier and consume more power for the rotating
motor when compared to single-point or matrix \gls{tof} sensors. Matrix \gls{tof} sensors are lighter and have
no moving parts but the range of \gls{cots} modules is more limited (usually less than 10m) when compared to
single-point \gls{tof} sensors \cite{dronelidar} due to power and safety constraints.

A hybrid between a \gls{tof} sensor and a \gls{lidar} sensor can be achieved by mounting one \gls{tof} sensor on
a gimbal, so that the sensor can be pointed at multiple directions. This allows for the measurement of multiple points
while keeping the long ranges of the single-point sensor.

\subsection{Radar}

Radar modules are increasingly used in drones for applications such as obstacle detection,
collision avoidance, and terrain mapping \cite{radaraltimetro2}. By emitting radio waves and
analyzing the reflected signals, radar sensors provide precise distance and velocity measurements,
even in low visibility conditions such as fog, rain, or darkness, where optical sensors may
struggle \cite{introtoradar}. The ability of radars to detect objects over longer distances and through
cluttered environments makes it ideal for autonomous navigation and enhancing flight safety.
In addition, advanced radar systems, such as millimeter-wave radars, are used for altimetry and
precise positioning during landings, particularly in drones operating in challenging outdoor environments. 

Although radar modules offer many advantages for drones, they also have notable disadvantages.
One key limitation is their high power consumption, which can reduce the drone's battery life,
particularly in smaller systems. Additionally, radar systems can struggle with resolution limitations,
especially for small or closely spaced objects, making them less effective for detecting fine details
compared to optical sensors like LIDAR or cameras. Environmental factors, such as electromagnetic
interference, can also affect performance and introduce noise into the radar signal. Furthermore,
cost can be a barrier, as advanced radar systems, such as millimeter-wave radars, tend to be more
expensive than simpler alternatives like ultrasonic or \gls{ir} sensors. Lastly, the complexity of
data processing and integration with other onboard systems requires advanced algorithms and additional
computational resources, which may complicate the drone's design.

\section{Target location}

Of all types of sensors analyzed in this work, optical ranging sensors are the mort capable and
precise considering weight, volume, price and energy consuption. Light-based sensors, such as
\gls{ir} or laser rangefinders, fit well within the technical constraints of the design. Radar
is too complex, bulky, and expensive. Ultrasonic sensors are simply not precise enough.

As the angle of the target relative to the drone module is known via the camera, the only
measurement required would be the distance from the target to the module. The rangefinder measures
this distance from above, and the geometry errors get corrected in firmware.

In order to keep costs low, use \gls{cots} components and maintain the flexibility of the module within
the constraints of the flight level specified by \gls{anac}, a hybrid \gls{tof} using a gimbal will be
tested alongside a system that simulates a drone using computer vision to test the feasibility
and accuracy of such a system.

